\section{Implementation}
\label{sec:implementation}

\subsection{Python package architecture}
\label{sec:package-architecture}

\bsi\ is distributed as a standard Python package (\texttt{src} layout,
\texttt{pyproject.toml}-based build) targeting Python 3.10 and above and
tested on CPython 3.10--3.13 across Linux, macOS, and Windows in continuous
integration (Section~\ref{sec:testing}). The installed package exposes the
layers of Figure~\ref{fig:architecture} as Python subpackages
(\texttt{bioseqinsight.core}, \texttt{.models}, \texttt{.services},
\texttt{.structures}, \texttt{.workflows}, \texttt{.io}, \texttt{.gui}) plus
a \texttt{cli} module. The \textbf{core computational package has no
required third-party runtime dependencies}: every nucleotide and protein
calculation in Section~\ref{sec:core-calculations} is implemented against
the standard library alone. Two optional extras are provided: \texttt{http}
installs the \texttt{requests} library for a connection-pooling HTTP
transport (the standard-library \texttt{urllib} transport is used
automatically when it is absent, with identical externally observable
behaviour), and \texttt{validation} installs Biopython
\citep{cock2009biopython}, which enables the cross-validation test module of
Section~\ref{sec:cross-validation} but is not required to run the software.
Installation is a single command:

\begin{lstlisting}
pip install bioseqinsight
\end{lstlisting}

with no compiler, no scientific numerical stack, and no system libraries
required for the core package. This is a deliberate design constraint
(Section~\ref{sec:architecture}) rather than an incidental property: the
target user includes biologists installing the software on a personal
laptop, for whom a missing compiler or an unsatisfiable binary dependency is
a common point of failure for scientific Python tools.

\subsection{Command-line interface}
\label{sec:cli}

Every capability of the graphical interface is reachable from the command
line, which is the primary interface used to generate the reproducible
figures and tables in Section~\ref{sec:evaluation} and is fully usable on a
headless machine or within a continuous-integration pipeline. Subcommands
cover nucleotide analysis (\texttt{dna}), protein analysis
(\texttt{protein}), structure retrieval and validation (\texttt{structure}),
batch processing of a FASTA file (\texttt{batch}), reproducible project
management (\texttt{project}), cache inspection (\texttt{cache}), effective
configuration (\texttt{config}), and version/environment reporting
(\texttt{version}). Exit codes distinguish success (0), user input error
(2), and runtime failure (3), so the interface composes directly into shell
scripts and continuous-integration checks. The verdict for each structure is a plain sentence in the report and in every
export. The three forms below are the message templates in the validation module,
shown for the exact, partial and mismatch cases:

\begin{lstlisting}
M4: Exact biological identity confirmed (identity 100.0%, coverage 100.0%);
    the structure cross-references <ACCESSION>.
partial_structure: Partial structure of the requested protein: the structure
    cross-references <ACCESSION> and matches the query at <x>% identity over
    the region it contains, but it covers only <y>% of your sequence. It is a
    real structure of this protein, not a full-length one.
M1 (measured mismatch): WARNING: the retrieved structure does not correspond
    to the requested protein (identity <x>%, coverage <y>%), which is below the
    threshold for even a related record. It is associated with <ACCESSION>.
\end{lstlisting}

\subsection{Graphical interface}
\label{sec:gui}

A four-tab Tkinter desktop interface (Sequence Analysis, Protein Analysis,
Structure Retrieval, Batch Analysis) calls the same underlying API as the
command line. Network operations execute on a background thread and post
results back to the Tkinter main loop through a queue, so a slow or stalled
external request does not freeze the interface; all widget mutation remains
confined to the main thread as Tkinter requires. The Structure Retrieval tab
displays, for every queried resource, its attempt history, whether it
recovered from a transient failure, whether the response was served from
cache, and a coloured M0--M4 badge with the identity and coverage measures
behind it (Table~\ref{tab:m0m4}); a retrieved structure can be opened in an
embedded 3Dmol.js \citep{rego20153dmol} viewer or handed to an external
viewer such as PyMOL \citep{pymol} or ChimeraX
\citep{pettersen2021chimerax} when one is installed. Tkinter ships with most
standard Python distributions; on platforms where it is absent, launching
the graphical interface prints an explanatory message and the equivalent
command-line invocation rather than raising an unhandled exception, so that
absence of a display toolkit degrades gracefully to the fully capable
command-line interface rather than making the software unusable.

The formatting logic that turns a result object into the text shown in the
graphical interface is implemented in a module with no Tkinter import
(\texttt{gui.formatting}), specifically so that the exact wording a user
reads -- including every warning that a retrieved structure is not a
confirmed match -- is covered by the same offline, headless test suite as
the rest of the software (Section~\ref{sec:testing}), rather than being
exercised only by manual interaction with the graphical interface.

\subsection{Caching}
\label{sec:caching}

Responses from every external resource are cached on disk under a
per-resource namespace, keyed by a hash of the query and relevant request
parameters, with a configurable time-to-live (seven days by default) and an
explicit refresh option. Each cached payload is stored alongside a SHA-256
checksum computed at write time; on read, the checksum is recomputed and
compared, and a mismatch -- whether from disk corruption or external
tampering -- causes the entry to be discarded and treated as a cache miss
rather than served. Cache statistics (hit rate, entry count, on-disk size)
are available through both interfaces. Caching is not merely a performance
optimisation in this software: it is what makes repeated re-analysis of the
same sequence set, including the live benchmark of
Section~\ref{sec:live-benchmark}, practical without repeatedly loading the
queried public resources.

\subsection{Batch processing}
\label{sec:batch}

A FASTA file of arbitrary size is processed as a batch: each record's
alphabet (nucleotide or protein) is detected automatically or may be
specified explicitly, local analysis proceeds according to the detected
type, and structure retrieval (when requested) is dispatched across a
configurable thread pool, since it is I/O-bound on external network
requests rather than CPU-bound. A single malformed or unanalysable record is
isolated to its own output row, carrying its specific error, rather than
aborting the run; the throughput and scaling behaviour of local (structure
retrieval excluded) batch processing is reported in
Section~\ref{sec:performance}.

\subsection{Export and project management}
\label{sec:export}

Results export to CSV, TSV, or JSON with a fixed, documented column set
(\texttt{io.export.BATCH\_COLUMNS}), and to a self-contained HTML report with
all styling inlined and no external script or stylesheet reference, so that
the report remains renderable indefinitely after the analysis that produced
it, independent of any external resource's continued availability. The
project mechanism of Section~\ref{sec:reproducibility} is exposed through
both interfaces: \texttt{bioseqinsight project create/info/export/import}
on the command line, and a \emph{Project} menu in the graphical interface
offering the same operations.
